\subsubsection{Complexity and Importance}\label{sec:bernstein-vazirani-complexity-and-importance}

The Bernstein-Vazirani algorithm is mainly important because it gives a very \textbf{clear example of quantum query advantage}. The computational problem itself is simple:
\begin{equation*}
    f_w(x) = w \cdot x \mod 2
\end{equation*}
Where $w \in \left\{0,1\right\}^n$ is unknown. The goal is to recover the complete hidden bitstring $w$.

\highspace
\begin{flushleft}
    \textcolor{Red2}{\faIcon{hourglass-half} \textbf{Classical Query Complexity}}
\end{flushleft}
A deterministic classical algorithm can \hl{recover one bit of $w$ per query.} Using the standard basis vectors:
\begin{equation*}
    e_1 = (1,0,0,\ldots,0), \quad e_2 = (0,1,0,\ldots,0), \quad \ldots, \quad e_n = (0,0,0,\ldots,1)
\end{equation*}
We obtain:
\begin{equation*}
    f_w(e_1) = w_1, \quad f_w(e_2) = w_2, \quad \ldots, \quad f_w(e_n) = w_n
\end{equation*}
Therefore, recovering all $n$ bits requires $n$ oracle queries. Hence the classical query complexity is:
\begin{equation*}
    Q_{\text{classical}} = n \rightarrow \mathcal{O}(n)
\end{equation*}
This is an exact deterministic result: after $n$ queries, the hidden bitstring is known with certainty.

\highspace
\begin{flushleft}
    \textcolor{Green3}{\faIcon{check-circle} \textbf{Quantum Query Complexity}}
\end{flushleft}
The quantum algorithm prepares a superposition of all possible inputs, \textbf{queries the oracle once}, and uses phase kickback and interference to recover the entire string. After the final Hadamard layer, the state is:
\begin{equation*}
    \ket{\psi_{\text{final}}} = \ket{w} \ket{-}
\end{equation*}
Where the first $n$ qubits are in the state $\ket{w}$, and the last qubit (ancilla) is in the state $\ket{-}$. Therefore, one measurement of the data register returns $w$ with certainty (probability 1). Thus, the quantum query complexity is:
\begin{equation*}
    Q_{\text{quantum}} = 1 \rightarrow \mathcal{O}(1)
\end{equation*}

\highspace
\begin{flushleft}
    \textcolor{Green3}{\faIcon{tachometer-alt} \textbf{What Kind of Speedup is This?}}
\end{flushleft}
\hl{The Bernstein-Vazirani algorithm gives an $n$-to$1$ reduction in oracle queries.} This is a real quantum advantage in the \textbf{query model}, but it is important to state it correctly.

\highspace
It is not an exponential speedup in the ordinary running-time sense. The quantum circuit still uses:
\begin{itemize}
    \item $n$ data qubits;
    \item $n+1$ Hadamard gates before the oracle;
    \item $n$ Hadamard gates after the oracle;
    \item Up to $n$ CNOTs inside the oracle;
    \item $n$ measurements.
\end{itemize}
Therefore, the \hl{total gate cost still grows with $n$.} The advantage is specifically that the expensive black-box function $U_f$ is queried only once. So the correct statement is that the \hl{Bernstein-Vazirani algorithm provides a query-complexity advantage}, \textbf{not a constant running time advantage}.

\highspace
\begin{flushleft}
    \textcolor{Green3}{\faIcon{link} \textbf{Relationship with Deutsch's Algorithm}}
\end{flushleft}
Bernstein-Vazirani is closely related to Deutsch's algorithm. Both algorithms use:
\begin{itemize}
    \item An ancilla initialized in \ket{1}.
    \item A Hadamard gate preparing the ancilla in \ket{-}.
    \item An oracle of the form:
    \begin{equation*}
        \ket{x}\ket{y} \mapsto \ket{x}\ket{y \oplus f(x)}
    \end{equation*}
    \item Phase kickback to encode the function value in the phase of the data register.
    \item A final Hadamard transformation to recover the hidden information.
    \item Deterministic measurement: the hidden information is recovered with probability 1.
\end{itemize}
The \textbf{main difference} is the \textbf{problem being solved}.
\begin{itemize}
    \item Deutsch's algorithm uses one data qubit and determines whether a\break Boolean function is constant or balanced.
    \item Bernstein-Vazirani uses $n$ data qubits and recovers an entire hidden bitstring:
    \begin{equation*}
        w \in \left\{0,1\right\}^n
    \end{equation*}
\end{itemize}
Thus, \hl{Bernstein-Vazirani is a multi-qubit extension of the phase-kickback idea used in Deutsch's algorithm.}

\highspace
\begin{flushleft}
    \textcolor{Green3}{\faIcon{link} \textbf{Relationship with Simon's Algorithm}}
\end{flushleft}
Bernstein-Vazirani and Simon's algorithm are both hidden-structure problems, but they recover different kinds of information.
\begin{itemize}
    \item \textbf{Bernstein-Vazirani}. The function is:
    \begin{equation*}
        f_w(x) = w \cdot x \mod 2
    \end{equation*}
    The hidden object is a bitstring $w$, and the algorithm returns it directly with \textbf{one query}: \ket{w}.


    \item \textbf{Simon's algorithm}. As we will see in the next section, the function satisfies the promise:
    \begin{equation*}
        f(x) = f(x \oplus a)
    \end{equation*}
    Where $a \neq 0$ is a hidden period. A single run does not return $a$. Instead, it returns a string $y$ satisfying:
    \begin{equation*}
        a \cdot y = 0 \mod 2
    \end{equation*}
    The \textbf{circuit must be repeated several times} to collect enough independent equations to reconstruct $a$.
\end{itemize}
Therefore, \hl{Bernstein-Vazirani returns the hidden string directly; Simon returns constraints on the hidden string.} Bernstein-Vazirani is simpler and deterministic after one run, while Simon introduces the more advanced idea of recovering hidden structure from repeated inference constraints.

\highspace
\begin{takeawaysbox}
    The Bernstein-Vazirani problem hides a bitstring:
    \begin{equation*}
        w \in \left\{0,1\right\}^n
    \end{equation*}
    Inside the Boolean function:
    \begin{equation*}
        f_w(x) = w \cdot x \mod 2
    \end{equation*}
    Classically, the bits of $w$ are recovered one at a time, requiring $n$ queries.

    \highspace
    The quantum algorithm prepares:
    \begin{equation*}
        \ket{0}^{\otimes n} \ket{1}
    \end{equation*}
    Applies Hadamard gates, and queries the oracle once. Because the ancilla is in \ket{-}, the oracle introduces the phase:
    \begin{equation*}
        \left(-1\right)^{w \cdot x}
    \end{equation*}
    The final Hadamard layer decodes this phase pattern:
    \begin{equation*}
        \ket{+} \mapsto \ket{0}, \quad \ket{-} \mapsto \ket{1}
    \end{equation*}
    The resulting state is:
    \begin{equation*}
        \ket{\psi_{\text{final}}} = \ket{w} \ket{-}
    \end{equation*}
    So measuring the data register returns the hidden bitstring $w$ with certainty (probability 1).
\end{takeawaysbox}